\section{Methodology}

% Plan: State the identification principle that organizes all experiments.
Our design follows one identification principle: change the linguistic route to
a belief while holding the world fixed, then change the identity of the
belief's competitors while holding the linguistic contrast fixed. The first
intervention establishes that form matters; the second reveals what that form
controls. We separate hypothesis formation from held-out prediction and
independent confirmation so that a descriptive sensitivity cannot promote
itself into a compositional law. The complete experimental inventory appears in
Appendix~\ref{app:evidence-inventory}.

\subsection{The answer is a typed competition}

% Plan: Define the latent world and explain why typed candidates identify source competition.
Each latent world specifies a queried agent, a comparison agent, an object, an
initial location, a true final location, and agent-indexed information states.
The event sequence is generated once and reused across all surface forms. Its
answer candidates are typed as the queried agent's belief, reality, the
comparison agent's state, or an object distractor. We construct role conflict,
object conflict, and joint role-plus-object conflict by changing which typed
competitors disagree with the instructed source. Candidate strings remain
invariant within a latent-world block. Consequently, a changing margin cannot
be attributed to a changing ground truth or answer vocabulary; it records
which source defeats which competitor.

% Plan: Make explicit the invariant checks that prevent surface changes from changing the task.
Every crossed block is audited for an identical event manifest, queried target,
candidate set, and instruction target. All 30 worlds in the same-item,
multi-realization, factorial, and first compound analyses pass these invariance
checks. The final confirmation uses 20 newly generated worlds with zero
normalized latent-signature overlap with the earlier pool. The rejected
cross-pool main-effect test likewise uses zero item overlap. A replay-equivalent
intermediate pool with complete identity overlap is retained only as an
execution check and is excluded from scientific replication claims.

\subsection{Descriptions change the route, not the world}

% Plan: Introduce the high-level linguistic interventions and their theoretical roles.
The first intervention family varies description order, reference form,
negation/evidential wording, and role labels. These contrasts manipulate which
source is introduced first, how explicitly its discourse referent is resolved,
and whether competing sources are excluded or the target evidence is
affirmatively licensed. The full inventory contains 10 controlled arms plus the
original description (Appendix~\ref{app:linguistic-arms}). Description order,
reference form, and negation/evidential wording are preregistered confirmatory
axes in the topology cross; role label versus name is retained as exploratory.

% Plan: Explain the realization and construction manipulations that separate robustness from accessibility.
We next re-express the same event structure as a canonical narrative,
structured event log, event-wise report, or compact indirect report. The
critical construction family crosses active versus genuine passive voice,
global versus event-local embedding, and reportative versus factive predicates.
These operations alter discourse prominence, grammatical argument realization,
embedding locality, or source commitment without changing the event sequence.
The compound test holds out all passive-by-embedding combinations while
estimating primitive voice and locality effects from the non-compound cells.

% Plan: Document exact-length and candidate-token controls.
We evaluate target context lengths of 1,024 and 2,048 tokens. Neutral padding is
inserted outside the answer-bearing event block, and exact length is checked
with the model tokenizer after the complete description and target block are
assembled. Every source decision reports full exact-length coverage. All
confirmatory answer candidates are required to be single tokens, preventing
multi-token normalization or stopping decisions from entering the margin.

\subsection{A margin identifies the winning source}

% Plan: Specify the model and inference procedure without implying training or hidden-state access.
Experiments use a fixed Qwen2.5-7B-Instruct checkpoint
\citep{qwen2024qwen25}. No weights are updated and no free-text decoding is
used for the primary evidence. The scorer obtains next-token logits at the
answer position and evaluates every typed candidate under the same prompt. The
checkpoint, tokenizer, data type, device map, and source manifests are fixed
before each confirmatory matrix; implementation details appear in
Appendix~\ref{app:reproducibility}.

% Plan: Define the answer-interface binding margin.
For prompt $x_i(c,T,a)$, let $\ell_i(y)$ be the next-token log probability of
candidate $y$. If $y_i^\star$ is the source licensed by the instruction and
$\mathcal{Y}_i^{-}$ is the set of typed competitors, the binding margin is
\begin{equation}
M_i(a,T,c)=\ell_i(y_i^\star)-\max_{y\in\mathcal{Y}_i^{-}}\ell_i(y).
\label{eq:binding-margin}
\end{equation}
A positive margin means that the instructed source wins the answer interface;
the magnitude records its distance from the strongest typed competitor. This
definition is intentionally behavioral: it does not assume a localized belief
representation or source-monitoring circuit.

% Plan: Define the paired linguistic effect and topology interaction.
For a linguistic contrast with arms $a^+$ and $a^-$, define
\begin{equation}
\Delta_i(T,c)=M_i(a^-,T,c)-M_i(a^+,T,c).
\label{eq:language-effect}
\end{equation}
Our primary topology estimand is the excess effect under joint conflict,
\begin{equation}
\Gamma_i(c)=\Delta_i(T_{\mathrm{role+object}},c)
-\tfrac{1}{2}\left[\Delta_i(T_{\mathrm{role}},c)
+\Delta_i(T_{\mathrm{object}},c)\right].
\label{eq:topology-interaction}
\end{equation}
This difference-in-differences removes the average single-conflict effect and
isolates how language form responds to concurrent typed competitors.

\subsection{Falsification before confirmation}

% Plan: Explain how each experiment changes the epistemic status of the claim.
Each transition in the argument is triggered by a failed explanation. Pooled
language effects are tested on an item-disjoint pressure pool before typed
conflict is introduced. The typed interaction is then challenged across four
surface realizations; its reportative reversal motivates the accessibility
hypothesis. Primitive voice and locality costs are estimated in a factorial and
frozen before unseen compounds are scored. Only after those compounds suggest a
saturating class do we register its residual sign and predictive controls and
test them on 20 new latent worlds. Later evidence therefore answers a new
question rather than adding rows to an earlier grid.

% Plan: State the negative controls and no-promotion rules.
Three controls prevent a descriptive pattern from being mistaken for
independent evidence.
First, the replay-equivalent pool is not counted as independent evidence.
Second, role label versus name is not promoted with the preregistered axes.
Third, a held-out additive predictor must beat the no-voice control before
additive composition can be claimed. It does not; the residual pattern is used
only to formulate the subsequent saturation hypothesis. Independent
confirmation then tests that hypothesis without selecting among saturation
operators.

\subsection{Uncertainty and confirmatory criteria}

% Plan: Define the cluster bootstrap and why the latent world is the sampling unit.
All estimands are first reduced within latent world by averaging the two context
lengths. We use nonparametric item-cluster bootstrap intervals, resampling
latent worlds while preserving every topology, construction, length, and
instruction arm belonging to a world. Unless stated otherwise, intervals are
95\% percentile intervals from 8,000 resamples with fixed analysis seeds. This
avoids treating correlated surface realizations as independent observations.

% Plan: State each confirmatory gate explicitly.
The topology claim requires the preregistered direction and a directional
confidence interval for all three primary axes, with the same direction at both
context lengths. Cross-realization confirmation requires an aggregate
directional interval and the predicted sign in at least three of four
realization families. Independent saturation requires a positive lower interval
bound in both reportative and factive families, together with improvements of
both weakest-link and average-saturation predictors over the frozen additive
null. Full cell-level results and decision criteria are reported in the
appendix.
